\chapter{Definitions}

Following key terms are used throughout this diploma project:\par\bigskip

\begin{description}[style=multiline, labelwidth=5cm, leftmargin=5.5cm]

  \item[Alert Fatigue]
       Cognitive overload experienced by Security Operations Center (SOC) analysts due to an excessive volume of low-value or false-positive security alerts generated by SIEM and IDS tools.

  \item[Asset Context]
       Operational metadata characterizing a monitored host or service within an organization's Configuration Management Database (CMDB), including business criticality, vulnerability exposure (EPSS), and network centrality.

  \item[Event Sequence Analysis]
       Temporal and stage-wise tracking of alerts targeting a specific asset across the phases of the MITRE ATT\&CK kill chain, capturing burst progression and multi-stage correlation.

  \item[Explainability by Design]
       An intrinsic property of a scoring model where the output priority is a transparent mathematical composition of domain-meaningful factors, requiring no post-hoc approximation (Faithfulness $= 1.0$).

  \item[Risk-Based Alerting (RBA)]
       A SIEM paradigm that modifies alert severity using static asset or user risk multipliers, evaluated in this work as a contextual baseline without multi-stage sequence memory.

\end{description}
